\def\ChapterCode{JUF}
\chapter
[\CO{[JUF]} Rechtliche Grundlagen (.at) -- Fallbeispiele]
{Rechtliche Grundlagen (.at) -- Fallbeispiele}
\label{chp:JUS}
% \AasRsN{RS~III}

\setcounter{tocdepth}{10}
\mtcsetdepth{minitoc}{3}
\ChapterIntro
{%
Dieses Kapitel behandelt die allgemeinen und
tätigkeitsspezifischen rechtlichen Grundlagen. 

Anhand von konkreten Beispielen vermittelt es die notwendige
Kompetenz und Problemlösungsstrategien um im Alltag
auftretende rechtliche Probleme und Fragestellungen zu
lösen.
}
{%
\begin{description}
  \item[Maintainer:] Christof Koller
  \item[Autoren:] Christof Koller et. al.
%   \item[Reviewer:] --
  \item[Version:] {\DocVersion} ({\DocVersionInfo})
  \item[SHA1:] {\DocGitCommit}
\end{description}

{\TxTeilkapitelAASS}
}
% \minitoc
\mtcsetdepth{minitoc}{\value{DefaultMinitocDepth}}
% \AuthNotes{
% Changelog:
% 2009-02-20: GABS: Fertig für Kurs 2009/02
% 
% 2009-02-25: EDIT Koller
% 
% 2009-04-22: GABS: Abkürzungen \enquote{Pat.}, \enquote{KA}, \enquote{KH}
% bereinigt und gegen Langform ausgetauscht, wie besprochen.}





\clearpage


\section{Ein langer Dienst}
\label{topic:langer-dienst}



{\centering
\EE{{\sffamily\Large Eine Geschichte aus 17 Rechtsfällen}}
\par}

\Text{Zweck}{%
Behandlung rechtlicher Fragen anhand von
Praxisbeispielen, die in einer zusammenhängenden Geschichte
erzählt werden. Dabei wird insbesondere versucht auf die
Besonderheiten des Rettungswesens Bezug zu nehmen und die
sich daraus ergebenden rechtlichen Probleme einfach
darzustellen. Die Darstellung beinhaltet aber nur einen
Überblick über die behandelten Rechtsgebiete sowie eine
Konkretisierung auf die am häufigsten gestellten
Rechtsfragen.
}{}{}{}{}



\Text{Aufbau}
{%
Jeder einzelne Rechtsfall wird in der Geschichte nach folgendem Aufbau
gestaltet:

\begin{enumerate}
    \item Sachverhalt: Schilderung der Geschichte
    \item Rechtsfrage: Erkennung des sich stellenden Problems
    \item Rechtslage: rechtliche Beurteilung des Problems
    \item Rechtsfolgen: drohende Konsequenzen und Ausweichmöglichkeiten
\end{enumerate}
}{}{}{}{}


\Text{Vorkommende Personen}
{%
\begin{center}
\begin{tabularx}{\linewidth}{XX}
\textbf{\textcolor{red}{A}}lex, Fahrer

\textbf{\textcolor{red}{S}}epp, Transportführer

\textbf{\textcolor{red}{B}}ernd, Zivildienstleistender

\textbf{\textcolor{red}{NA}}tascha, Notärztin

&
Herr \E{Hauzu}, 1. Patient

Frau \E{Uninformiert}, 2. Patientin

Frau \E{Finito}, 3. Patientin

Herr \E{Schmerz}, 4. Patient

Frau \E{Aufgeregt}, Angehörige

Dr. \E{Wasbringst}, Aufnahmearzt

Frau \E{Hilflos}, 5. Patientin

\\
\end{tabularx}
\end{center}
}{}{}{}{}

\Text{Der Tag beginnt \ldots}
{%
Unsere Geschichte beginnt an einem warmen
Frühjahrs\-morgen auf den die drei Sanitäter Alex, Sepp und
Bernd schon lange gehofft hatten, da sie endlich nach dem
langen Winter einen angenehm warmen Sonnentag genießen
wollen. Doch wie immer an einem solchen Freudentag, ruft die
Arbeit mit einem ganz normalen Sanitätsdienst und seinen
typischen Problemen.
}{}{}{}{}


\subsection{Gefahrenbekämpfung}

\subsubsection{Fall 1. Randalierender Patient}
\label{topic:unterbringung-jus}
\label{topic:jus-langer-dienst-fall-1}

\TODO{GABS}{yyyy-mm-dd}{Subsubsection: Layout-Anpassung notwendig.}{Dieser
Unterunterabschnitt entspricht nicht dem Standard-Layout!}


\Text{Sachverhalt}{%
Die erste Ausfahrt beginnt natürlich direkt nach
Dienstbeginn, ohne dass wenigstens ein Kaffee getrunken
wurde. Der auf unsere RTW-Mannschaft wartende
Patient, Herr Hauzu, hingegen hat die ganze Nacht
durchgemacht und ist schon weit\-läufig bekannt für seine
tobend aggressiven Wutausbrüche im Zuge seiner
psychiatrischen Erkrankung. Als Alex, Sepp und Bernd
eintreffen, droht er -- aus einiger Entfernung -- sich oder
einen anderen zu verletzen, falls man sich ihm auch nur
nähern würde.
}


\subparagraph{Variante \enquote{unmittelbare Bedrohung}:} Herr Hauzu
stürmt auf Alex zu, schmeißt ihn zu Boden und holt zum
Schlag aus.

\Text{Rechtsfrage}
{Darf man sich gegen Hrn. Hauzu wehren? Muss Herr Hauzu
eingewiesen werden?} 


\Text{Rechtslage}{%
Für die erste Frage sind die Selbsthilferechte maßgebend,
vgl. \FullPrettyRef{topic:selbsthilferechte}:

\begin{itemize}
  \item Notwehr: Abwehr eines Angriffes
  auf ein eigenes oder fremdes Rechtsgut
  
  \Speech{\color{DefaultA}Solange Herr Hauzu nur droht, kann keine
  Maßnahme in Notwehr ausgeführt werden, da kein
  unmittelbarer Angriff vorliegt. Greift Herr Hauzu jedoch
  tatsächlich an, liegen die Voraussetzungen der Notwehr vor
  und es können entsprechende Abwehrhandlungen getroffen
  werden.}
  \item \E{Notwehrüberschreitung im Affekt}:
  \Speech{\color{DefaultA}Sollten Alex, Sepp und Bernd
  fälschlich einen Angriff angenommen haben, sind sie
  entschuldigt, wenn sie aus Furcht gehandelt haben und eine
  maßgerechte Person sich gleich verhalten hätte.}
  \item  \E{Anhalterecht Privater}:
  
  \Speech{\color{DefaultA}Solange keine gerichtlich strafbare Handlung
  (oder ein entsprechender Verdacht) vorliegt (z.\,B.
  eine Körperverletzung oder Sachbeschädigung durch
  einen Angriff), kann das
  Anhalterecht Privater nicht angewendet werden.}
\end{itemize}

Bezüglich einer \enquote{Einweisung} ist das
\I[Unterbringungsgesetz!Fallbeispiel]{Unterbringungsgesetz}
(\FullPrettyRef{topic:unterbringungsgesetz}) zu
berücksichtigen:
\begin{itemize}
  \item \Speech{\color{DefaultA}Da Herr Hauzu eine
  psychiatrische Erkrankung hat, durch diese eine ernstliche
  Gefährdung für andere droht
  und keine verhältnismäßigeren Maßnahmen möglich
  erscheinen, ist eine Unterbringung einzuleiten
  (Verständigung der Sicherheitspolizei und in weiterer
  Folge des Amtsarztes).}
\end{itemize}


}{}{}{}{}

\Text{Rechtsfolge}{%
\FazitLike{
Alex, Sepp und Bernd haben somit \E{grundsätzlich} das
Recht, eventuelle Gewaltausbrüche des Herr Hauzu gegen sich
und andere abzuwehren. In der Situation, in der bloß aus einiger
Entfernung mit Gewalt gedroht wird, \E{ohne dass eine
unmittelbare Gefahr vorliegt}, ist
es \E{verhältnismäßig}, zurück zu treten, die Wohnung zu
verlassen und die Polizei zu verständigen. 

\subparagraph{Variante \enquote{unmittelbare Bedrohung}:} Alex,
Sepp und Bernd können das verhältnismäßig gelindeste
Zwangsmittel wählen, um sich gegen den unmittelbaren Angriff
zu \E{schützen}. Sie sind jedenfalls \E{verpflichtet},
unverzüglich die Polizei und über diese den Amtsarzt zu
verständigen. Trotzdem sollte Gewalt durch Sanitäter immer
als letzte Option verwendet werden, auch wegen dem Risiko
sich durch das Anhalten selbst zu verletzen. Daher der
Grundsatz: Selbstschutz geht vor Fremdschutz (bzw.
übertriebenen Heldenmut).
}
}{}{}{}{}


\subsubsection{Fall 2. Gefahr durch Hund}
\label{topic:jus-langer-dienst-fall-2}

\Text{Sachverhalt}{%
Zwischenzeitlich befreit sich der \enquote{kleine} (65\Kg* schwere)
Rottweiler-Kampfhund von Herrn Hauzu aus einem Nebenzimmer,
erblickt die drei \enquote{Eindringlinge} und möchte mit
fletschenden Zähnen sein Revier verteidigen.
}
{}{}{}{}



\Text{Rechtsfrage}{Was darf gegen den Hund getan
werden?\bigskip}{}{}{}{}


\Text{Rechtslage}{%
Grundsätzlich sind Tiere in der Rechtsordnung als Sachen
anzusehen. Für die Selbsthilferechte bedeutet dies,
dass kein \E{menschlicher} Angriff, sondern eine
Gefahrensituation vorliegt. Es stellt sich nun die Frage
nach dem Notstand:

\Speech{\color{DefaultA}Da der Angriff unmittelbar drohend
ist, und auch die sonstigen Voraussetzungen der
Notstandssituiationen vorliegen, dürften sich die Sanitäter
mit verhältnismäßigen Mitteln gegen den Hund zur Wehr setzen
(\I[Notstand!rechtfertigender!Fallbeispiel]{rechtfertigender
Notstand}). Dies wäre auch dann möglich, wenn der Angriff
nur subjektiv, aber nachvollziehbar angenommen wird
(\I[Notstand!entschuldigender!Fallbeispiel]{entschuldigender
Notstand}).
Sinvollerweise werden die Sanitäter jedoch den Rückzug
antreten (möglichst ohne das Tier zu provozieren) und
\E{keine} Notstandshandlung durchführen. }
}{}{}{}{}

\Text{Rechtsfolge}{%
\FazitLike{
Alex, Sepp und Bernd dürfen sich mit den schonendsten Mitteln
verteidigen, um ihr eigenes Leben oder ihre Gesundheit zu
schützen. Unabhängig davon treten sie jedoch den Rückzug
an.}
}{}{}{}{}


\Cave{Schon vor Betreten der Wohnung sollte ein
Hund von seinem Besitzer in einen  separaten Raum
weggesperrt werden.}


\subsection{Heilbehandlung}

\subsubsection{Fall 3. Blutzuckermessung =  Heilbehandlung}
\label{topic:jus-langer-dienst-fall-3}
\index{Heilbehandlung}

\Text{Sachverhalt}{%
Durch die rechtzeitig eintreffende Polizei, kann die Gefahr
durch den Hund beseitigt und Herr Hauzu festgehaltenen
werden. Trotz der Zuweisung durch den etwas später
eingetroffenen Amtsarzt, möchte Bernd, gewissenhaft wie er
ist, einen vollen Neuro-Check mit einem Blutzuckertest
durchführen.

}{}{}{}{}



\Text{Rechtsfrage}{%
Wann liegt eine  \I{Heilbehandlung} vor? Darf der Blutzuckertest
durchgeführt werden?
}{}{}{}{}


\Text{Rechtslage}{%


\Speech{\color{DefaultA}Sowohl die geplante neurologische
Untersuchung, als auch die Blutzuckermessung stellen eine
Heilbehandlung dar (\FullPrettyRef{topic:heilbehandlung}).
Aufgrund der krankheitsbedingt fehlenden Einsichts- und
Urteilsfähigkeit muss von einer Aufklärung bzw. Einwilligung
des Patienten (in seinem Interesse!) abgesehen
werden.}\ZFootnote{(Die Messung des Blutzucker mit einer
Lanzette war früher umstritten, da manche
Rettungsorganisationen davon ausgingen, dass es sich um
einen nicht dem Sanitäter erlaubten Eingriff am menschlichen
Körper handelt. Dieser Streit wurde durch eine Novellierung
des SanG gelöst, indem man die Blutzuckermessung für den
Sanitäter ausdrücklich zuließ.)}
}{}{}{}{}

\Text{Rechtsfolge}{%
\FazitLike{Da ein kompletter Neuro-Check bei einem
neurologisch auffälligen Patienten medizinisch indiziert und
gesetzlich erlaubt ist, darf jeder der drei Sanitäter eine
solche Maßnahme lege artis durchführen. Eine Aufklärung ist
aufgrund der mangelnden Einsichts- und Urteilsfähigkeit
nicht möglich.}
}{}{}{}{}

\subsection{Selbstbestimmungsrecht}

\subsubsection{Fall 4. Aufklärung und Einwilligung}
\label{topic:jus-fall-aufklaerung}
\label{topic:jus-langer-dienst-fall-4}
\index{Aufklärung}
\index{Einwilligung}



\Text{Sachverhalt}{%
Nachdem Herr Hauzu gebändigt, zur Gänze untersucht und
untergebracht wurde, freut sich die RTW-Mannschaft schon auf
den wohlverdienten Kaffee am Stützpunkt. Doch wie es kommen
musste, können Alex, Sepp und Bernd nicht einrücken, da sie
einen Folgetransport zugewiesen bekommen. Die vorgefundene
Patientin Frau Uninformiert ist eine in der 30. Woche
schwangere Mutter, die erst letzte Woche ihren 18.
Geburtstag feierte. Grund für ihren Anruf sind starke
Kopfschmerzen und leichte Seh- und Höreinschränkungen. Bei
der weiteren Untersuchung durch Sepp werden geschwollene
Knöchel, ein Blutdruck von \RRmmHg{140}{90} und nach
beiläufiger Aussage der Patientin starker Harndrang mit
süßlichem Geruch diagnostiziert.

Schnell ist Sepp klar, dass es sich um eine \I{EPH-Gestose}
(\ref{topic:eph-gestose}) handelt, die zur Eklampsie führen
kann und einen raschen, aber möglichst reizarmen Transport
auf die Gynäkologie benötigt. Da er aber wenig von jungen
Müttern hält, erachtet er es auch nicht für notwendig Frau
Uninformiert über die nächsten Geschehnisse zu informieren,
sondern ordnet nur streng Befehle an, die von Frau
Uninformiert \enquote{gefälligst} durchzuführen sind.
}{}{}{}{}

\Text{Rechtsfrage}{%
Besteht eine Aufklärungspflicht und was umfasst sie? Muss eine
Einwilligung erfolgen und wenn ja, in welcher Form?
}{}{}{}{}


\Text{Rechtslage}{%
\Speech{\color{DefaultA}\FullPrettyRef{topic:aufklaerung}:
Die Patientin Frau
Uninformiert ist volljährig und zeigt keine Kennzeichen
einer fehlenden Einsichts- und Urteilsfähigkeit. Ihre
Entscheidungsbefugnis ist gegeben,
daher ist sie vollinhaltlich aufzuklären.}

}{}{}{}{}




\Text{Rechtsfolgen}{%
\FazitLike{%
Aufgrund der unterlassenen Aufklärung machte sich Sepp strafbar nach
folgenden Bestimmungen:

\begin{itemize}
  \item \EE{Strafrecht}:
  \begin{itemize}
    \item Eigenmächtige  Heilbehandlung  ({\S}\,110 StGB):
    \begin{itemize}
      \item nur bei \enquote{Gefahr in
      Verzug}\marginlabel{Gefahr in Verzug: \E{unmittelbar}
      drohende oder gegenwärtige Gefahr, die ein Abwarten
      nicht zulässt bzw.
      unmittelbares Handeln verlangt} benötigt es keine
      Aufklärung bzw. Einwilligung und es dürfen alle nach
      Ansicht der Einsatzkräfte direkt notwendigen Maßnahmen
      gesetzt werden
    \end{itemize}
    \item Körperverletzung \Z{({\S}{\S}\,83 ff StGB)}:
    \begin{itemize}
      \item keine Strafbarkeit, solange die Heilbehandlung
      medizinische indiziert und lege artis
      durchgeführt wurde
    \end{itemize}
  \end{itemize}
  \item \EE{Zivilrecht}:
  \begin{itemize}
    \item
    Schadenersatz wegen Körperverletzung \Z{({\S}{\S}\,1293
    ff ABGB)}:
    \begin{itemize}
      \item {für das Zivilrecht ist die Einwilligung
      aufgrund einer Aufklärung notwendige Voraussetzung für
      die rechtfertigende Heilbehandlung ${\rightarrow}$
      mangelt es an ihr, kann Schadenersatz verlangt werden}
    \end{itemize}
  \end{itemize}
\end{itemize}
}
}{}{}{}{}

\subsubsection{Fall 5. Patientenverfügung}
\label{topic:patientenverfuegung}
\label{topic:jus-langer-dienst-fall-5}


\Text{Sachverhalt}{%
Nach Ablieferung der Frau Uninformiert drängt Alex seine
Kollegen zu einer schnellen Abfahrt, um endlich seinen Kaffee
zu bekommen. Doch da macht ihm die Leitstelle schon wieder
einen Strich durch die Rechnung, als sie das Team zu einem
weiteren BO beruft.

Dieser liegt in einem kleinen Pflegeheim, in dem eine
aufgeregte Heimpflegerin auf Alex, Sepp und Bernd zustürmt
und sie zu der 70 jährigen, regungslosen Patientin, Frau
Finito, bringt.
Gleichzeitig hält die Heimpflegerin den Sanitätern eine
Patientenverfügung vor, die Frau Finito vor über 5 Jahren auf
einer Serviette unterschrieben hat und nach der \enquote{bei
ausbleibenden Kreislauf keine Wiederbelebungsmaßnahmen
gesetzt werden} dürfen. Deshalb hatte die Pflegerin es bis
jetzt auch unterlassen mit einer Reanimation zu beginnen. Da
sie aber ebenso nicht weiß, ob die Verfügung bindend ist,
rief sie die Rettung und setzte sich wartend neben die
Patientin.
}{}{}{}{}

\Text{Rechtsfrage}{%
Was ist eine Patientenverfügung? Ab wann und wie weit ist
sie bindend?
}{}{}{}{}

\Text{Rechtslage}{%
Siehe \FullPrettyRef{topic:patientenverfuegung}.

\Speech{\color{DefaultA}Die \enquote{Verfügung} erfüllt
nicht die strengen
Kriterien einer verbindlichen Patientenverfügung: }

\begin{itemize}
  \item[\Irritant] Schriflichkeit vor Notar bzw.
  Rechtsanwalt
  \item[\Irritant] Aufklärung durch Arzt
  \item[\Irritant] 5-jährige Gültigkeit
\end{itemize}

\Speech{\color{DefaultA}Auch die Beachtlichkeit ist aufgrund der
Beiläufugkeit der Abfassung (auf Serviette) zu verneinen.}


 }{}{}{}{}

\Text{Rechtsfolge}{%
\FazitLike{
Die Patientenverfügung gilt mit über 5 Jahren, der fehlenden
ärztlichen Aufklärung und dem mangelnden Beiziehen eines
Rechtsverständigen nicht mehr als verbindlich und wäre somit
als bloß beachtlich einzustufen.


Doch aufgrund des zu allgemeinen Inhalts, der Beiläufigkeit der
Verfügung (Serviette als Unterlage) und wegen des hohen Alters der
Verfügung kann kein konkreter und aktueller Wille der
Patientin erkannt werden, was dazu führt, dass das Wohl
der Patientin das oberste Gebot bleibt.

Mit einer Reanimation ist zu beginnen.

}
}{}{}{}{}

\subsection{Rechte und Pflichten des Sanitäters}

\subsubsection{Fall 6. Unterlassung der Notarztnachforderung}
\label{topic:jus-langer-dienst-fall-6}



\Text{Sachverhalt}{%
Endlich am Stützpunkt angekommen, ist Alex auch schon der
Erste der sich zum Kaffeeautomaten begibt. Entkräftet durch
die einsetzende \enquote{Hypokoffeinämie} kramt er nach
Kleingeld und kaum als er es einwurfbereit in Händen hält,
tönt die nächste Ausfahrt durch die Lautsprecher mit dem
Hinweis einer dringenden Berufung. Alex geht aber in seiner
typisch ruhigen Manie, glücklich durch die
enteral\footnote{Enteral: den Darm betreffend.} erlangte
Sedierung, zum Fahrzeug, mit dem Spruch auf den Lippen:
\enquote{immerhin konn des
Wagl{\textquoteright} eh net ohne mie foahrn!}}{}{}{}{}



\begin{wrapfigure}{r}{5cm}
\includegraphics[width=\linewidth]{./images/teaser/nef-passat-kontur.jpg}
\end{wrapfigure}Nach nur kurzer Anfahrtszeit begegnet Alex,
Sepp und Bernd im ersten Stock eines kleinen Reihenhaus der
55 jährige, adipöse Herr Schmerz mit Schmerzen hinter der
Brust, die in den linken Arm ausstrahlen und ihm selbst in
Ruhelage mit Atemnot und kalten Schweißausbrüchen plagen.
Für Sepp ist sofort klar, dass es sich um einen Herzinfarkt
handelt.
Er möchte schleunigst einen NEF nachfordern, doch Alex
meint, dass er keine Lust hätte schon wieder auf einen
\enquote{Boder\footnote{"Bader": früher Kundiger der
Heilbäder, heute umgangssprachlich für Arzt}} zu warten und
der Patient auch selbst runter zum Fahrzeug gehen sollte, da
er es im Kreuz hätte und sowieso noch eine ca. 20 Minuten
lange Anfahrtszeit zur Krankenanstalt durch den schon
einsetzende Morgenstau bewältigen müsste.

\Text{Rechtsfrage}{%
Besteht die Pflicht für Sanitäter zur Nachforderung eines Notarztes?}
{}{}{}{}

\Text{Rechtslage}{%
\II[Hilfeleistung!unterlassene]{Unterlassene Hilfeleistung} - {\S}\,95 StGB

\begin{itemize}
    \item 
    jeder Durchschnittsmensch ist zur Durchführung von
    \E{ihm zumutbaren Hilfeleistungen} verpflichtet. Die
    Zumutbarkeit bestimmt sich dabei nach dem individuellen
    Stand an Ausbildung, Fähigkeiten und Kenntnissen. Aber
    schon zu den Grundlagen einer Ersten Hilfe, die von
    jedem Menschen mit Besuch eines Führerscheinkurses
    erwartet werden kann, gehört das Absetzen eines Notrufes
    zur  Anforderung von geeigneten Rettungskräften.
\end{itemize}

\EE{{\S}\,4 SanG}

\begin{itemize}
    
    \item 
    jeder Sanitäter, unabhängig von seinem Ausbildungsstand,
    ist zur Durchführung von lebensrettenden Sofortmaßnahmen
    verpflichtet, was gleichzeitig beinhaltet, dass er bei
    Notwendigkeit eine unverzügliche Anforderung des
    Notarztes oder sonstigen Arztes zu veranlassen hat.
    
\end{itemize}

\TODO{GABS}{2010-05-28}{Unterlassene Hilfeleistung: Abgrenzung von
Unterlassung von schädlichen Handlungen.}{
In der Diskussion wird oft jede Unterlassung mit unterlassener Hilfeleistung
gleichgesetzt. Es sollte ein gegenteiliger Hinweis erfolgen, zB.

\enquote{Anmerkung: Die Unterlassung von schädlichen Handlungen, d.\,h. Handlungen die
\E{nicht} dem Patientenwohl dienen, stellt \E{keine} unterlassene Hilfeleistung
dar. Schädliche Dinge zu unterlassen ist genauso wichtig wie nützliche Dinge zu
tun.} 

Die Beurteilung, welche Maßnahme schädlich ist, ist oft stark vom Einzelfall
abhängig. }


}{}{}{}{}

\Text{Rechtsfolge}{%
\FazitLike{
Würde Alex die zweckmäßige Nachforderung eines Notarztes unterlassen,
macht er sich gleichzeitig nach strafrechtlichen ({\S}\,95 StGB) und
verwaltungsrechtlichen Bestimmungen ({\S}\,9 in
Verbindung mit {\S}\,53 SanG) strafbar. }
}{}{}{}{}

% \clearpage
\subsubsection{Fall 7. Sanitäterkompetenzen}
\label{topic:kompetenzen-sanitaeter}
\label{topic:jus-langer-dienst-fall-7}



\Layout{60}{Sachverhalt}{%
Nach langem hin und her können Sepp und Bernd gemeinsam den
Alex überreden doch einen Notarzt nachzufordern und Sepp,
der gerade frisch gebackener NFS geworden ist, sieht endlich
seine Möglichkeit umfangreich an Herr Schmerz zu werken.
Dabei möchte er neben der Verabreichung von
Nitrolingual-Spray auch gleich einen venösen Zugang legen,
da er das sowieso als ehemaliger Fleischhauer besser als
jeder
Arzt kann.}
{}
{./images/venflon_in_situ.jpg}
{Ein blutiger Venflon vom Fleischhauer?}
{Gabriel.}



\Text{Rechtsfrage}
{Was sind die \T{Kompetenzstufen} der einzelnen
Sanitäter-Ausbildungsstufen (RS, NFS, NKA, NKV, NKI)? Darf
man die Kompetenzgrenzen (bei \E{Gefahr in
Verzug}) überschreiten?}{}{}{}{}



\Text{Rechtslage}{%
Die \EE{Kompetenzstufen} sind im
SanG\ZFootnote{{\S}{\S}\,\VonBis{9}{12} SanG} geregelt.
\FullPrettyRef{tab:sang-kompetenzstufen}, gibt eine
Übersicht.


\subparagraph{Überschreiten der Kompetenzstufen}
Die Durchführung von Maßnahmen außerhalb des eigenen
Kompetenzstandes, führt zu einer Verwaltungsstrafe nach SanG
bzw.
einer gerichtlichen Strafe wegen Körperverletzung und
zivilrechtlichen Schadenersatzanspruch. Selbst bei Wissen
über Gefahren, Risiken und Durchführungsart einer
kompetenzmäßig höheren Maßnahme, gibt es keine
Rechtfertigung durch {\S}\,95 StGB (Unterlassene
Hilfeleistung).

Dies bedeutet weiters, dass einfache Maßnahmen, die zwar das
Berufsgesetz nicht aufzählt, aber von jedem Laien (mitunter nach
kurzer Einweisung) durchgeführt werden dürfen, auch von
Sanitätern gesetzt werden dürfen bzw. iSd. {\S}\,95 StGB gesetzt
werden müssen.
}{}{}{}{}

\Text{Rechtsfolge}{%
\FazitLike{%
Sepp darf keinen venösen Zugang legen oder sogar Medikamente darüber
verabreichen, selbst  wenn Gefahr in Verzug vorliegt. }
}{}{}{}{}


\subsubsection{Fall 8. Rechtswidrige Anordnung durch Notarzt}
\label{topic:jus-langer-dienst-fall-8}

\Layout{60}{Sachverhalt}{%
Sepp konnte sich nun dazu überreden lassen keinen Venflon zu
legen und möchte wenigstens Nitro verabreichen. Gerade noch
rechtzeitig bemerkt er eine offene Packung von Viagra am
Tisch liegen. Herr Schmerz meint, als er Sepp die Packung
betrachten sieht, dass das \enquote{Zeug} nichts nutzt, da
er vor zwei Stunden eine Tablette eingenommen hätte und noch
immer \enquote{nichts steht}. Kurze Zeit später trifft der
NEF mit der Notärztin Natascha ein, die den Patienten
oberflächlich untersucht und eine sofortige Nitro-Verabreichung
zur Bekämpfung des hohen Blutdrucks (\RRmmHg{200}{110})
anordnet, obwohl sie von Sepp über die gefundene
Viagra-Packung informiert wurde.
}{

}
{./images/teaser/emergency_physician_vehicle_00800.jpg}
{}
{}


\Text{Rechtsfrage}{Ist der Sanitäter verpflichtet eine
rechtswidrige bzw. gefährliche Anordnung des Notarztes (oder
eines anderen Vorgesetzten) zu folgen?}
{}{}{}{}



\Text{Rechtslage}{%
\EE{Ablehnung von rechtswidrigen Anordnungen}:

\begin{itemize}
  \item
  Zivildienstleistende: {\S}\,22 Abs. 2 Zivildienstgesetz:
  \begin{itemize}
    \item
    \enquote{Er darf die Befolgung einer Weisung nur dann ablehnen,
    wenn [{\dots}] die Befolgung gegen strafgesetzliche Vorschriften
    verstoßen würde.}
  \end{itemize}
  \item
  Haupt- und Ehrenamtliche:
  \begin{itemize}
    \item
    {\S}\,4 Abs. 1 S. 2 SanG
    \item
    {\S}\,7 i.\,V.\,m. {\S}\,879 Abs. 1 ABGB
    \item
    {\S}\,22 Abs. 2 ZDG analog
  \end{itemize}
\end{itemize}

Aufgrund der im SanG fehlenden konkreten Bestimmung eines
Ablehnungsrechtes von Sanitätern, kann ein solches nur kompliziert
gewonnen werden aus:

\begin{itemize}
    \item allgemeinen Verpflichtung zur Wahrung des Patientenwohls nach {\S}\,4
    SanG    
    \item natürlichen Rechtsgrundsätzen iSd. {\S}\,7 ABGB und dem
    Nichtigkeitsgebot bei Sittenwidrigkeit des {\S}\,879 ABGB    
    \item analoger Anwendung des Zivildienstgesetz (ZDG) 
    \item Schutz der Grundrechte gem. Art.~2 EMRK bei
    Lebensgefahr und Art.~3 EGC bei sonstigen Verletzungen
    der Unversehrtheit.    
\end{itemize}
}{}{}{}{}

\Text{Rechtsfolge}{%

\FazitLike{%
Die RTW-Mannschaft darf daher die für den Patienten	schädigende Anordnung ablehnen. Natascha hingegen macht
sich strafbar wegen Körperverletzung und riskiert uU.
zivilrechtliche Schadenersatzansprüche
(${\rightarrow}$ Nitro-Gabe nicht lege artis).}
}{}{}{}{}


\subsubsection{Fall 9. Dokumentationspflicht}
\index{Dokumentationspflicht}
\label{topic:jus-langer-dienst-fall-9}
\label{topic:dokumentationspflicht-fallbeispiel}


\Text{Sachverhalt}
{Genervt von der Auseinandersetzung der Notärztin mit den
Sanitätern, beginnt Herr Schmerz wild vor Schmerzen die
Beteiligten anzubrüllen, dass sie doch endlich ihre
\enquote{gottverdammte Arbeit} tun sollen.
Erschrocken von dem verbalen Angriff des Herrn Schmerz tritt
Bernd einen Schritt zurück und hört ein dumpfes Klirren
unter seinem Fuß. Als er hinunterblickt sieht er, dass unter
seinem Schuh Porzellanscherben verstreut sind. Auch Herr
Schmerz bemerkt die Scherben und wird noch wütender, denn
die Scherben stammen von einer (angeblich) wertvollen
Blumenvase, die am Boden abgestellt war und nun von Bernd
zerstört wurde. Bernd erinnert sich aber, die Scherben schon
beim Betreten des Raumes gesehen zu haben, da er extra über
sie hinweggestiegen ist.

Wochen später beschwert sich Herr Schmerz beim
Arbeitgeber von Alex, Sepp und Bernd über das
\enquote{randalierende Verhalten} von Bernd. Bernd seinerseits ist
überzeugt, dass er sich korrekt verhalten hat, doch er kann sich so recht an den alten Fall nicht mehr
erinnern, da er zwischenzeitlich unzählige weitere
Einsätze hatte und der Vorfall nicht dokumentiert
wurde.}{}{}{}{}



\Text{Rechtsfrage}{%
Besteht eine \I{Dokumentationspflicht} und wenn ja, in welchem Umfang?
}{}{}{}{}

\Text{Rechtslage}{
\Speech{\color{DefaultA}Für
Sanitäter besteht die Pflicht, nicht nur gesetzte medizinische Maßnahmen zu dokumentieren,
sondern auch andere relevante Umstände. Daher sind
die Behauptungen des Patienten in die Dokumentation
aufzunehmen. \FullPrettyRef{topic:dokumentationspflicht}}

}{}{}{}{}~

\Text{Rechtsfolge}{%
\FazitLike{Bernd hat keine ausreichende Dokumentation durchgeführt und kann sich
somit nicht gegen die Beschuldigungen zur Wehr setzen.}
}{}{}{}{}

\subsubsection{Fall 10. Verschwiegenheits- und Auskunftspflicht}
\label{topic:jus-langer-dienst-fall-10}




\Text{Sachverhalt}
{Durch die Aufregung verschlechtert sich allmählich der
Zustand des Herr Schmerz und das RTW-Team samt NEF-Mannschaft
sind sich einig, dass ein schleunigster Abtransport
durchzuführen ist. Noch bevor sie den Patienten
umlagern können erscheint Frau Aufgeregt, die Angehörige des
Herr Schmerz, und verlangt über alles informiert zu werden.
}{}{}{}{}

\Text{Rechtsfrage}{Was umfasst die
Verschwiegenheitspflicht der Sanitäter und Ärzte? Besteht
auch eine Auskunftspflicht derselben?
}{}{}{}{}



\Text{Rechtslage}{%
\Speech{\color{DefaultA}Jeden Sanitäter trifft die
grundsätzliche Pflicht zur Verschwiegenheit. Nur in
Ausnahmesituationen darf oder muss Auskunft über Geheimnisse
gegeben werden.(\FullPrettyRef{topic:verschwiegenheitspflicht},
\FullPrettyRef{topic:auskunftspflicht})}
}{}{}{}{}

\Text{Rechtsfolge}{%
\FazitLike{%
Frau Aufgeregt zählt als bloße Angehörige nicht zum Kreis der
Auskunftsberechtigten, solange der Patient die Weitergabe von
Informationen nicht gestattet. }
}{}{}{}{}


\subsubsection{Fall 11. Fortbildungspflicht}
\label{topic:jus-langer-dienst-fall-11}
\label{topic:jus-fall-fortbildung}


\Text{Sachverhalt}{Nachdem Frau Aufgeregt
beruhigt wurde, wird Herr Schmerz auf den Sessel umgelagert,
um über die enge Treppe den ersten Stock zu verlassen. Doch
handelt es sich bei dem Tragsessel um einen neuen
Stryker{\texttrademark}-Sessel, der dem Bernd noch
vollkommen neu ist, da es den Sessel zum Zeitpunkt seiner
RS-Ausbildung noch gar nicht gegeben hat. Aber genau wegen
dieser Neuheit hatte sein Arbeitgeber eine interne
Fortbildung über die \enquote{Neuen Medizinprodukte}
abgehalten, zu der Bernd aber nicht gekommen ist, da er
meinte sowieso schon alles zu wissen und zu keinen
Fortbildungen verpflichtet ist.

Sepp hat hingegen als vorbildlicher Praxisanleiter eine
Einschulung nach MPG erhalten.}{}{}{}{}



\Text{Rechtsfrage}
{Besteht eine Fortbildungspflicht?\bigskip}
{}{}{}{}



\Text{Rechtslage}{%
\Speech{\color{DefaultA}Nach dem Medizinproduktegesetz
(MPG) hat jeder Anwender eines Medizinproduktes vor der
Anwendung eine gerätespezifische Einschulung zu absolvieren.
Davon unabhängig besteht eine fachbezogene
Fortbildungspflicht sowie eine
Rezertifizierungspflicht.}

\noindent\FullPrettyRef{topic:fortbildungspflicht},
\FullPrettyRef{topic:rezertifizierungspflicht} 

}{

}{}{}{}

\Text{Rechtsfolge}{%
\FazitLike{%
Die für den vorliegenden Fall entscheidende Fortbildung gem.
MPG wurde von Bernd nicht besucht, was ihm nicht nur eine
Verwaltungsstrafe, sondern auch eine gerichtliche Strafe und
zivilrechtlichen Schadensersatzanspruch einbringen kann.
Denn sollte es im Umgang mit dem Medizinprodukt zu
Körperverletzungen kommen, wird aufgrund der mangelnden
Materialkenntnis die Behandlung als nicht-lege-artis
gewertet. Um einer solchen Strafe zu entgehen genügt es,
dass Bernd noch vor Ort, aber vor Anwendung, durch eine
\enquote{geeignete Person} eingeschult wird.

Als \enquote{geeignete Person} gelten vorrangig
Medizinproduktebeauftragter,  Lehrsanitäter oder
Praxisanleiter sowie jene, die von der Rettungsorganisation
als geeignet genannt werden.

Da Sepp Praxisanleiter ist, darf er vorab Bernd einweisen.
}
}{}{}{}{}

\subsection{Schadenersatz}
\label{topic:schadenersatz}

\subsubsection{Fall 12. Verletzung des Patienten durch
Sanitäter:}
\label{topic:jus-langer-dienst-fall-12}



\Text{Sachverhalt}
{Nachdem nun Sepp unseren Bernd kurzum erklärt hat, wie er
mit den Stryker{\texttrademark}-Sessel umzugehen hat,
unterlässt er trotzdem das Anschnallen des Patienten, da er
meint, dass \enquote{eh kana däs mocht}. Leider ist Herr
Schmerz kein Fliegengewicht und der Treppenaufgang 
derart verwinkelt, dass es alle Kraft benötigt, um den
Tragsessel ruhig zu halten.

Fast am Ende des Abgangs angekommen flitzt auf einmal die
aufgeschreckte Katze des Herr Schmerz an den beiden
\enquote{Trägern} vorbei, wodurch der vorangehende Bernd
seine Konzentration verliert, ausrutscht und samt dem
Patient seitlich nach vorne auf einen Kasten und die
darunter sich versteckende Katze namens
\emph{\enquote{Kitty}} %\emph{\enquote{Muschi}} 
fällt. Bernd kann sich beim Sturz wenigsten selber abstützen
und bleibt (auch durch die abfedernde Wirkung der Katze)
unverletzt. Doch Herr Schmerz prallt mit Schädel und
Brustkorb ungeschützt gegen den Kasten und wird unter dem
hinter ihm nachkommenden Tragsessel begraben. Er zieht sich
neben einer RQW auf
der Stirn mehrere gebrochene Rippen zu.
}

\clearpage 
\Text{Rechtsfrage - Schadenersatz wegen
Körperverletzung}{ Wann wird man schadenersatzpflichtig? Was
umfasst der Schadenersatz? Wer wird zur Haftung gezogen?
}
{}{}{}{}



\TextSlide{Rechtslage}{~

\EE{Rechtliche Grundlage eines
Schadenersatzanspruches}:
   Die Grundlagen ergeben sich aus dem
    Bürgerlichen Recht\ZFootnote{Grundlagen des
    Schadensersatzanspruches: Allgemein nach
    {\S}{\S}\,1293\,ff ABGB und speziell für den Fall der
    Körperverletzung nach {\S}{\S}\,\VonBis{1325}{1327} ABGB}.


\EE{Voraussetzungen der Schadenersatzpflicht}, vgl.
\FullPrettyRef{topic:haftbarkeitsvoraussetzung}:

\begin{enumerate}
    \item \EE{Schaden}: Nachteil am Körper, an der
    Gesundheit: 
    
    \Speech{\color{DefaultA}RQW Stirn, gebrochene Rippen}
    \item \EE{Kausalität}: bei Wegdenken der
    Schädigerhandlung, bleibt auch der Erfolg aus:
    
    \Speech{\color{DefaultA}Bei Anschnallen des Patienten
    hätte er mit an Sicherheit grenzender Wahrscheinlichkeit
    einen weitaus geringeren Schaden erlitten.}
    \item \EE{Rechtswidrigkeit}: objektiver Verstoß gegen
    Gesetze oder Vertragspflichten
    
    \Speech{\color{DefaultA}Das Anschnallen des Patienten ist laut
    Hersteller des Tragestuhls verpflichtend. Weiters
    existiert eine interne Dienstanweisung der Einrichtung,
    nach welcher Patienten nur angeschnallt in dem
    Tragstuhl berfördert werden dürfen.}
    \item \EE{Schuld}: subjektiver Verstoß gegen jene
    Sorgfaltspflichten, die man einhalten hätte können
    
    \Speech{\color{DefaultA}Aufgrund der erfolgten Geräteeinschulung
    gem. Medizinproduktegesetz muss den handelnden
    Personen subjektiv klar sein, dass ein Anschnallen
    objektiv geboten ist.}
\end{enumerate}

\EE{Umfang des Schadenersatzes bei
Körperverletzungen}:

\begin{itemize}
	\item Heilungskosten
	\item Verdienstentgang
	\item Schmerzengeld
\end{itemize}


\EE{Wer haftet wie}:

Schädigt der Sanitäter im Zuge seines Dienstes einen Patient
oder dessen Angehörigen, kann der Geschädigte nach
\Z{{\S}\,1313a} ABGB direkt gegen die Rettungsorganisation
seine Schadenersatzansprüche geltend machen: Ein Sanitäter
gilt in dieser Beziehung als Erfüllungsgehilfe der
Rettungsorganisation, wodurch er als verlängerter Arm der
Organisation deren vertragliche Pflichten (sorgfältige
Behandlung) erfüllt. Daher kann der Geschädigte direkt gegen
die Organisation alle Ansprüche wie gegen den Sanitäter
selber durchsetzen.
}{
\begin{itemize}
  \item Grundlage: Bürgerliches Recht (ABGB)
  \item Voraussetzungen:
  \begin{enumerate}
    \item {Schaden}
    \item {Kausalität}
    \item {Rechtswidrigkeit}
    \item {Schuld}
\end{enumerate}
\item Umfang
\begin{itemize}
	\item Heilungskosten
	\item Verdienstentgang
	\item Schmerzengeld
\end{itemize}
\end{itemize}
}
{}{}{}

\Text{Rechtsfolge}{%
\FazitLike{%
Im vorliegenden Fall ist der Schaden durch RQW und Frakturen
neben der Kausalität gegeben. Weiters verstößt es gegen die
gesetzlichen und vertraglichen Pflichten den Patient ohne
anzuschnallen zu transportieren, weshalb eine
Rechtswidrigkeit bejaht werden kann.

Auch hätten sich unsere Sanitäter anders, nämlich den
Patient anzuschnallen, verhalten können, wodurch ebenso die
Schuld gegeben ist. Daher kann Herr Schmerz
Schadenersatzansprüche gegen die Sanitäter oder,
wirtschaftlich vernünftiger, gegen die Rettungsorganisation
geltend machen.
}
}{}{}{}{}


\subsubsection{Fall 13. Sachbeschädigung}
\label{topic:jus-langer-dienst-fall-13}


\Text{Sachverhalt}
{Leider sind die Verletzungen am Patient nicht die einzigen
Schäden die entstanden sind, denn der umgeschmissene Kasten
war eine Biedermeier-Rarität mit einem Wert von stolzen
10\,000,-- Euro. Auch die Katze verwirkte als Auffangkissen
für Bernd ihr Leben und hatte, neben dem Sachwert von 350,--
Euro, auch einen besonderen emotionalen Wert für
Hrn. Schmerz.}{}{}{}{}



\paragraph{Rechtsfrage: Schadenersatz wegen Sachbeschädigung}

Wann wird man schadenersatzpflichtig? Was umfasst der
Schadenersatz? Wer wird zur Haftung gezogen?

\TextSlide{Rechtslage}{%
~

\EE{Rechtliche Grundlage eines
Schadenersatzanspruches}:
   Die Grundlagen ergeben sich aus dem
    Bürgerlichen Recht\ZFootnote{Grundlagen des
    Schadensersatzanspruches: Allgemein nach
    {\S}{\S}\,1293\,ff ABGB und speziell für den Fall der
    Sachbeschädigung {\S}\,1331 ABGB}.
% \begin{inparaitem}
%     \item Allgemein: \Z{{\S}{\S}\,1293 ff ABGB}; 
%     \item Konkret für Sachschäden: \Z{{\S}\,1331 ABGB}
% \end{inparaitem}

\EE{Voraussetzungen der Schadenersatzpflicht}:
Es gelten die gleichen Voraussetzungen wie bei Körperverletzungen.

\EE{Umfang des Schadenersatzes bei Sachbeschädigung}:
Maßgebend ist der 
\begin{itemize}
  \item \E{reale Schaden} (der
  tatsächlich eingetretene Schaden),
  \item ein \E{entgangener Gewinn} (ein in Zukunft drohender
  Schaden),
  \item sowie ein \E{ideeller Schaden}.
\end{itemize}

\bigskip


}{
\begin{itemize}
  \item Rechtliche Grundlage: Bürgerliches Recht (ABGB)
  \item Voraussetzung Schadenersatzpflicht:
  \begin{itemize}
    \item Wie bei Körperverletzungen
  \end{itemize}
  \item Umfang Schadenersatz bei Sachbeschädigung:
\begin{itemize}
    \item Realer Schaden (= tatsächlich eingetretener Schaden)
    \item Entgangener Gewinn (= in Zukunft drohender Schaden)
    \item Ideeller Schaden
  \end{itemize}
\end{itemize}
}{}{}{}

\Text{Rechtsfolge}{%
\FazitLike{Gleiche Lösung wie bei Fall 12.}
}{}{}{}{}

\subsection{Behandlungspflicht}


\subsubsection{Fall 14. Behandlungspflicht von Krankenanstalten, Ärzten,
Sanitätern und Durchschnittsbürgern:}
\label{topic:jus-langer-dienst-fall-14}
\label{topic:behandlungspflichtJUS}


\Text{Sachverhalt}{%
Vor Schmerzen und Trauer, um seine geliebte Katze
\E{\enquote{Kitty}},
%\E{\enquote{Muschi}},
beschimpft und verflucht Herr
Schmerz die Anwesenden mit wilden Gestikulierungen. Als er
dann wenigstens von Alex und Sepp vorbildlich versorgt wird,
beruhigt sich Herr Schmerz und lässt sich sogar mühelos ins
Auto transportieren.

Nach ausreichender Versorgung des Herr Schmerz begibt sich
der RTW auf die Reise zur Krankenanstalt. Doch ist leider
das einzig abbuchbare Bett, welches neben einer Internen-
auch eine Unfall-Ambulanz aufzuweisen hat, im
Wilhelminenspital, welches durch den großen Morgenstau gute
15 Minuten entfernt ist. Schon nach kurzer Fahrzeit klagt
Herr Schmerz über stärker werdende Atemprobleme und ein
rasselndes, brodelndes Atemgeräusch tönt durch den gesamten
Patientenraum. Es hat sich ein kardiales Lungenödem gebildet
und Natascha fühlt sich, trotz Verabreichung von Medikamenten,
nicht mehr in der Lage den Patient während der Fahrt
ausreichend zu behandeln.
Alex soll so schnell wie möglich das nächste Spital
anfahren. Durch eine kurze Abweichung der Route, gelangt das
Team zu einem näher gelegenen, geeigneten Spital und eilt
mit dem Patient zur nicht vorinformierten Aufnahme.

Der dortige Aufnahmearzt Dr. Wasbringst hat schon einen 22
Stunden Dienst hinter sich und ist alles andere als erfreut
über das \enquote*{Mitbringsel} des Teams. Wutentbrannt
brüllt er Natascha an, was sie sich eigentliche erlaube
trotz gesperrten Bettenkontingent seine Station anzufahren
und verweigert ihr
die Aufnahme des Herr Schmerz.}{}{}{}{}



\Text{Rechtsfrage}{%
Besteht eine Behandlungspflicht für Krankenanstalten? Besteht eine
Behandlungspflicht für medizinische Berufe? Besteht eine
ersthelferische Versorgungspflicht für den Durchschnittsbürger?
}{}{}{}{}


\TextSlide{Rechtslage}{%
\Speech{\color{DefaultA}Nicht nur Gesundheitsberufe wie
Ärzte oder Sanitäter haben eine Behandlungspflicht gegenüber
dem Patienten, sondern auch jeder Durchschnittsbürger.
Ebenso sind Krankenanstalten verpflichtet, akut
behandlungspflichtige Patienten aufzunehmen und ihnen
zumindest eine ärztliche Grundversorgung zukommen
zu lassen.
}
}{

}{}{}{}    


\Text{Rechtsfolge}{%
\FazitLike{%
Auch wenn das Bettenkontingent der Krankenanstalt
aufgebraucht ist, muss Dr.
Wasbringst den Patient aufnehmen und zumindest soweit
stabilisieren, dass er ohne Gefahren in das eigentliche
Zielspital weitertransportiert
werden kann.}
}{}{}{}{}




\subsection{Meldepflichtenrecht}

\subsubsection{Fall 15. Ansteckende Krankheiten}
\label{topic:jus-langer-dienst-fall-15}



\Layout{60}{Sachverhalt}{%
Als Herr Schmerz nach langem hin und her auf der
Abteilung für \E{Innere Medizin} aufgenommen wurde, erholen
sich Alex, Sepp und Bernd bei einem kurzen Plausch mit der
NEF-Mannschaft. Die Ruhe hält aber nicht lang an, als auch
schon der nächste Einsatz durchgegeben wird, der RTW und
NEF zum selben BO schickt.

Dort angekommen finden sie in einer Lacke von Blut die
25-jährige Frau Hilflos liegen, die sich selber
(querverlaufend) die Pulsadern aufgeschnitten hat. Als Sepp
beginnt sie zu versorgen, erzählt Frau
Hilflos, dass ihr Leben \enquote{sowieso keinen Sinn mehr
hat}, da sie ihr Freund mit
\index{HIV!Meldepflicht}\Abk{HIV} angesteckt hat.
}
{}
{./images/teaser/blut-tatort-jacke.jpg}
{}
{GABS.}





\Text{Rechtsfrage}{
Bei welchen Erkrankungen besteht eine Meldepflicht? In welchem
Ausmaß besteht eine solche Meldepflicht?
}{}{}{}{}



\Text{Rechtslage}{%
\Speech{\color{DefaultA}Bei einer Infektion mit HIV, ohne
das ein AIDS mit sichtbaren Krankheitszeichen vorliegt, hat
keine Meldung an die Behörde zu ergehen. Ebenso hat niemand
das Recht, aufgrund der HIV-Infektion Zwang auf den
Patienten auszuüben.}

}{}{}{}{}

\Text{Rechtsfolge}{%
\FazitLike{%
Nun ist zwar unser RTW-Team nicht zur Meldung von HIV an
eine Behörde verpflichtet, doch sehr wohl zur Meldung an die
weiterbehandelnden Gesundheitsberufe aufgrund der schon
erwähnten Auskunftspflicht gem.
\Z{{\S}\,7} SanG.}
}{}{}{}{}


\subsubsection{Fall 16. Strafbare Handlungen}
\label{topic:jus-langer-dienst-fall-16}


\label{topic:strafbare-handlungen}


\Layout{60}{Sachverhalt}{%
Die Schnittwunden am Handgelenk sind nicht die einzigen Verletzungen der
Frau Hilflos. 

Mehrere Hämatome\footnote{\emph{Hämatom:}
Bluterguss.} überziehen den Körper und ein großes Veilchen
ziert das rechte Auge der jungen Patientin. Als Natascha auf
diese weiteren Verletzungen eingeht, beginnt Frau Hilflos
weinend zu berichten, dass ihr Freund sie öfters geschlagen
hat und ihr mit dem Tod gedroht hat, falls sie die
Polizei verständigen würde.
}
{}
{./images/trauma/haematom-fehlpunktion-001.jpg}
{Hämatom.}
{GABS.}



\Text{Rechtsfrage}{%
Besteht eine Pflicht zur Anzeige von strafbaren Handlungen?
Wenn ja, für wen besteht diese Anzeigepflicht?
}{}{}{}{}



\Text{Rechtslage}{%
\Speech{\color{DefaultA}Es besteht ein strafrechtlicher
Sachverhalt, welcher eine ärztliche Anzeigepflicht für die
Notärztin und ein allgemeines Anzeigerecht für die
Sanitäter begründet. }

}{}{}{}{}

\Text{Rechtsfolge}{
\FazitLike{%
Natascha muss aufgrund ihrer gesetzlichen Anzeigepflicht die
Polizei nachfordern, die zur Einbringung einer Anzeige den
Sachverhalt aufnehmen wird. Dies muss sogar dann geschehen,
wenn die
Patientin selber keine Anzeige einbringen möchte.}
}{}{}{}{}




\subsection{Straßenordnung}
\label{topic:strassenverkehrsordnung}

\subsubsection{Fall 17. Blaulichtverwendung}
\label{topic:jus-langer-dienst-fall-17}


\Layout{60}{Sachverhalt}{%
Nach Durchführung der Behandlung und Anzeige gegen den
aggressiven Freund, begibt sich das RTW-Team erschöpft zum
Stützpunkt zurück.
Hungrig und müde quält sich die RTW-Mannschaft durch das
alltägliche Wiener Verkehrschaos, als Sepp auf einmal meint:
\Speech{\enquote{Geh, schalt doch das Hörndl und die Blitzer
ei, damit dee uns endli Plotz machen und wir einrücken
kännan!}}
}
{}
{./images/gabriel-sebastian-aass/IMG_1896_00800.jpg}
{}
{GABS.}



\Text{Rechtsfrage}{%
Wann darf das Blaulicht samt Folgetonhorn verwendet werden? Wie darf
bzw. muss man sich bei einer Blaulichtfahrt verhalten?
}{}{}{}{}



\Text{Rechtslage}{\Speech{\color{DefaultA}
Voraussetzung für eine Einsatzfahrt ist Gefahr im Verzug,
diese liegt im gegenwärtigen Fall nicht vor
(\FullPrettyRef{topic:einsatzfahrt}).}
}{}{}{}{}

\Text{Rechtsfolge}{%
\FazitLike{%
Würde unser RTW-Team tatsächlich Blaulicht und Folgetonhorn
einschalten, nur um schneller einrücken zu können, würden
sie sich für die falsche Blaulichtfahrt gerichtlich strafbar
machen. Bei einem -- auch  unverschuldeten -- Unfall müssten
unsere drei Sanitäter die
volle Haftung tragen.}
}{}{}{}{}

\vspace{1em}
Der weitere Tag verläuft ruhig mit den üblichen Patienten
bis zum \ldots


\begin{center}
\EE{{\sffamily\Large Feierabend.}}
\end{center}


\ChapterEnd
